\chapter{Limitations, Threats to Validity, and Future Work}
\label{ch:limitations}

\section{Internal validity}

The project addresses several common leakage routes, but residual threats remain:

\begin{itemize}
  \item The same public participants were examined across multiple experiment
  generations. Nested CV protects each run's inner fitting, but it cannot erase
  researcher-level adaptation after earlier outer outcomes were seen.
  \item V1 compares ranks and families without nested selection. Its maximum is
  exploratory.
  \item V2 calibration uses 24 windows per training person for computational
  efficiency. This is reproducible but may underrepresent some within-person
  variability.
  \item V3/V4 spectral and waveform branches cover different effective bands:
  spectra use 1--30 Hz, while waveform moments retain approximately 0.5--45 Hz.
  The comparison is between complete pipelines, not a perfectly isolated
  representation ablation.
  \item Repeated CV folds are correlated. The three repeat scores are too few for
  a reliable normal-theory confidence interval.
\end{itemize}

\section{Construct validity}

Balanced accuracy measures classification, not biomarker validity. The project
does not establish that any factor corresponds to a neurophysiological disease
mechanism. The original proposal's MMSE correlations, ANOVA of factor loadings,
topographic biomarker maps, dPLV connectivity, and Riemannian baselines remain
unimplemented. Because channel order was wrong in V1, its graph/electrode
interpretations must not be used.

The waveform moment is exact for mean squared linear delay-filter energy, but it
cannot represent arbitrary nonlinear waveform dynamics or long-range order.
Calling it ``raw training'' without this qualification would overstate the model.

\section{External validity}

The cohort is single-site, limited to 88 people, and includes only 23 \FTD{}
participants. Recording device, clinical workflow, derivative preprocessing,
population characteristics, and label criteria may differ elsewhere. No result
is calibrated for clinical probabilities, prospective use, or medical decisions.

\section{Statistical power and uncertainty}

The main unit is the person, not the window. This leaves modest test folds of
roughly 17--18 people and creates sensitivity to which \FTD{} cases are held out.
The 59.7\% mean is only about 5.2 percentage points above the V4 primary endpoint
and is chosen as a best family after comparison. A formal inferential claim would
need an external test cohort or a prespecified locked model, plus uncertainty
quantification appropriate for participant-level clustered data.

\section{Prioritized next experiments}

Future work should be sequenced rather than expanded into an unlimited search.

\begin{enumerate}
  \item \textbf{Lock an external-validation protocol.} Freeze one representation,
  one classifier, all preprocessing thresholds, and a single primary metric
  before obtaining a new cohort.
  \item \textbf{Acquire independent FTD participants.} New people are more valuable
  than duplicating existing windows. Multi-site recruitment would also test
  device/site generalization.
  \item \textbf{Run targeted ablations.} Frequency-match waveform and spectral
  branches; compare TT16, spectral-32, and fusion-32 under identical preprocessing
  and one fixed classifier.
  \item \textbf{Add the missing course baselines.} Implement tangent-space
  Riemannian covariance features, then compare them under the same nested splits.
  \item \textbf{Investigate connectivity carefully.} Implement dPLV or another
  phase-connectivity tensor with leakage-safe estimation and surrogate/null
  checks. Avoid presenting volume-conduction effects as disease connectivity.
  \item \textbf{Stability and interpretability.} Evaluate Tucker/TT subspace
  stability across folds, bootstrap electrode/frequency loadings by participant,
  and correct for multiple comparisons before claiming biomarkers.
  \item \textbf{Clinically separate objectives.} If MMSE association is studied,
  treat it as a distinct outcome with its own validation rather than adding MMSE
  to an EEG diagnosis classifier.
\end{enumerate}

\section{What should not be done next}

The project should not repeatedly expand the same 88-person hyperparameter grid
until 65--70\% appears. It should not split by window, oversample \FTD{} windows
as if they were new patients, merge the leftover directory destructively, or
report secondary family maxima as the primary endpoint. Those actions would
increase apparent performance more readily than scientific credibility.

